\chapter{3D Spatial Data \& Point Clouds}
\label{ch:3d_spatial_data}

Visual recognition tasks such as image classification, object detection, and instance segmentation have seen superhuman performance when handling 2D images. However, the real 3D world is considerably more complex than 2D projections. There is an intrinsic ambiguity in the $3D \rightarrow 2D$ mapping: after projection, different depths cannot be distinguished in the image plane, meaning some information is fundamentally lost. Multiple views of the same scene help resolve these ambiguities; stereo vision systems leverage calibrated cameras and point correspondences to enable 3D reconstruction. Finding correspondences relies on extracting and normalizing region content, computing local descriptors, matching these descriptors, and performing robust model fitting.

\section{3D Data Representation}

Unlike images, which inherently come in an array format, 3D data are encoded in various ways, each with unique properties requiring different models for visual recognition.

\subsection{Point Clouds}
Point clouds represent the canonical form of raw 3D data. They can be formalized as a set of points in $\mathbb{R}^d$. Formally, a point cloud $P$ is defined as $P \subset \mathbb{R}^d$. When $d=3$, it simply records the 3D coordinates $(x, y, z)$ of points. When $d>3$, additional information such as color (RGB triplets), surface normals, or data-driven features are associated with the spatial coordinates. Point clouds are highly popular in autonomous driving due to their direct output from sensors like LiDAR. 
\begin{insightbox}[Advantages and Disadvantages of Point Clouds]
\textbf{Pros}: Point clouds offer a highly compact representation by not allocating memory for empty regions. They can intrinsically adapt to provide denser representations in regions with finer details.
\textbf{Cons}: Points are scattered and unordered; they do not reside on a regular grid. Standard Neural Networks and CNNs are built for ordered vectors or grid-based arrays. Because points theoretically have zero volume, they do not explicitly represent the continuous 3D surface. Furthermore, the coordinate representation heavily depends on the chosen reference coordinate system.
\end{insightbox}

\subsection{Voxels}
Voxels represent the 3D world as a discrete uniform grid $N \times N \times N$. In its simplest form, it contains binary values denoting occupancy: $V \in \{0,1\}^{N \times N \times N}$. Voxels can be generated by discretizing point clouds—setting a cell to 1 if it contains points, and 0 otherwise. Intensity information or other features can also be stored inside voxels (e.g., in Medical Imaging like MRI and fMRI).
While conceptually simple and structurally akin to 2D images (making it easy to adapt 3D CNNs), voxel grids suffer from severe scalability issues. Volumetric data are typically exceedingly sparse (often $<5\%$ occupancy at moderate resolutions). Thus, capturing fine structures requires high spatial resolution, making the memory footprint computationally prohibitive.

\subsection{RGB-D and 2.5D Images}
RGB-D images add a fourth depth channel to the standard RGB pixel grid. These are considered \textit{2.5D images} because they do not capture the full 3D geometry; they suffer from occlusions (the structure behind visible objects is not captured) and missing values due to shadows or non-reflecting surfaces. However, their matrix structure makes neighborhood computations trivial using standard pixel adjacency.

\subsection{Implicit Functions and Meshes}
An implicit function models a 3D object as a function $f: \mathbb{R}^3 \rightarrow \{0,1\}$ (occupancy) or $f: \mathbb{R}^3 \rightarrow \mathbb{R}$ (Signed Distance Function, SDF). The object surface is extracted as the 0-level set of $f$. 
Meshes are the standard representation in Computer Graphics, represented as a collection $(P, E, F)$ where $P \subset \mathbb{R}^3$ are the vertices, $E \subset \mathbb{R}^3 \times \mathbb{R}^3$ are the edges connecting vertices, and $F \subset \mathbb{R}^3 \times \mathbb{R}^3 \times \mathbb{R}^3$ are the triangular faces. Meshes are adaptive and efficient for flat surfaces, allowing native interpolation, but feeding them into Neural Networks requires specialized Graph CNN architectures.

\section{3D Imaging Sensors and Datasets}

\subsection{Sensors}
\begin{itemize}
    \item \textbf{Computed Tomography (CT) Scan}: A diagnostic imaging technique that measures multiple X-rays from different angles. Tomographic reconstruction algorithms (like filtered back projection) produce virtual slices of a body, discretized into a grid of voxels containing the CT number (density).
    \item \textbf{LiDAR (Light Detection and Ranging)}: Active sensors comprising a laser source (IR, UV) and a sensor measuring reflected light time-of-flight (ToL). They emit pulses to sweep the scene and construct precise 3D point clouds. However, LiDAR lacks color information and is sensitive to rain and weather conditions. Examples include Innovusion Falcon K and Ouster OS1.
    \item \textbf{RADAR}: Operates similarly to LiDAR but uses electromagnetic emissions from an antenna. It is robust and allows detecting moving objects and estimating their relative velocity. RADAR point clouds are typically much sparser than LiDAR.
    \item \textbf{Coded Light}: Uses synchronized projectors and cameras. The projector emits high-frequency black-and-white patterns. The signature serves as a descriptor for dense correspondence matching.
\end{itemize}

\subsection{Datasets}
\begin{itemize}
    \item \textbf{ModelNet}: The first large-scale 3D CAD dataset containing 151,128 models over 660 categories. Often used in 10 or 40 class variants.
    \item \textbf{ShapeNet}: A richly annotated dataset of 53,100 CAD models across 55 categories. It includes rendered images and ShapeNetPart for segmentation.
    \item \textbf{ScanNet}: Scans from 1,513 indoor scenes with 3D camera poses, surface reconstructions, and semantic voxel labeling.
    \item \textbf{Outdoor Datasets}: KITTI 360 and Waymo Open Dataset offer LiDAR data labeled with 3D bounding boxes and semantic segmentation.
    \item \textbf{BraTS}: Brain tumor segmentation dataset containing 369 multimodal MRI scans with spatial resolution $240 \times 240 \times 155$.
\end{itemize}

\section{Deep Learning on Voxelized Data}

A conventional 2D CNN utilizes locality-aware processing through convolutional layers followed by dense MLP layers. For 3D volumetric data (or videos treating time as the Z-axis), we extend this to 3D Convolutions. A 3D convolutional layer computes its output activation $a$ for a given spatial location and channel $k$ as:
\begin{equation}
a(x, y, z, 1) = \sum_{(u, v, s) \in U, k} w^1(u, v, s, k) X(r+u, c+v, d+s, k) + b^1
\end{equation}
The channel dimension aggregates all input features without spatial processing. To achieve spatial invariance globally, Global Average Pooling (GAP) or Global Max Pooling (GMP) is employed:
\begin{equation}
F_k = \frac{1}{H \cdot W \cdot Z} \sum_{(x,y,z)} f_k(x, y, z) \quad \text{for } k=1,\dots,C
\end{equation}

\begin{insightbox}[3D CNN Enhancements]
Standard 3D CNNs suffer from massive parameter counts and high memory requirements. They can be improved by:
1. \textbf{mlpconv layers} to reduce overfitting risk.
2. \textbf{Auxiliary subvolume supervision} to predict object classes from portions of the latent space.
3. \textbf{Anisotropic Probing}: using elongated 3D kernels to increase the receptive field.
4. \textbf{Orientation-Pooling}: training a backbone on standard orientations, generating features for multiple rotated 3D versions, and pooling them before final classification.
\end{insightbox}

To circumvent the staggering computational cost of native 3D convolutions, alternative representations like \textbf{Multi-View CNNs} render 3D shapes from multiple virtual cameras (e.g., 30 views) and apply 2D CNNs, followed by view-pooling. Another approach for video data is \textbf{(2+1)D Convolution}, which factors a 3D convolution into a 2D spatial convolution followed by a 1D temporal convolution, adding non-linearity and improving optimization. Additionally, \textbf{OctTrees} leverage adaptive cell sizes to losslessly reduce memory consumption in sparse voxel grids, although their kernels are rigidly constrained (e.g., $3^3 = 27$ voxels).

\section{Single Image Depth Estimation}

Estimating a depth map $M = (x,y,z)$ from a single RGB image is inherently ill-posed due to the \textit{scale-depth ambiguity}: a small close object perfectly mimics a large distant object. Yet humans infer depth using visual cues like perspective, texture gradient, occlusions, light/shadows, and defocus. Modern approaches use Fully Convolutional Networks (like an Encoder-Decoder with skip connections) to predict depth maps in a supervised manner.

A naïve $\ell_2$ pixel-wise loss $D(y, y^*) = \|y - y^*\|_2^2$ fails because global scale errors dominate. Thus, a \textbf{Scale Invariant Loss} is utilized to prioritize relative depth:
\begin{equation}
D(y, y^*) = \frac{1}{n} \sum_{i=1}^n \left( \log y_i - \log y^*_i + \alpha(y, y^*) \right)^2
\end{equation}
By setting $d_i = \log y_i - \log y^*_i$, the loss decomposes algebraically:
\begin{equation}
D(y, y^*) = \frac{1}{n} \sum_i d_i^2 - \frac{\lambda}{n^2} \left( \sum_i d_i \right)^2
\end{equation}
where $\lambda$ controls the scale invariance (e.g., $\lambda=0.5$). 
To preserve local geometric structures and sharp boundaries, a \textbf{Loss on Gradients} is introduced:
\begin{equation}
l_{grad} = \frac{1}{n} \sum \left[ (\nabla_x d_i)^2 + (\nabla_y d_i)^2 \right]
\end{equation}
Furthermore, a \textbf{Loss on Normals} evaluates the cross-product of tangent vectors $l_{norm} = \cos(n, n^*)$.
Despite these advanced losses, single-image depth estimators suffer from epistemic uncertainty on Out-of-Distribution (OOD) objects and struggle to generalize reliably. Position in the frame often acts as a stronger cue than apparent size.

\section{Deep Learning on Point Clouds}

Due to the unordered, continuous, and unstructured nature of point clouds, voxelization implies data loss and inefficiency. The seminal \textbf{PointNet} processes raw point clouds by employing symmetric functions to guarantee permutation invariance. 
According to the \textbf{Deep Sets} theorem, a function operating on a set $X$ is invariant to permutations if it can be decomposed as $\rho \left( \sum_{x \in X} \phi(x) \right)$, where $\phi$ and $\rho$ are approximated via MLPs.

\subsection{PointNet and PointNet++}
PointNet uses an MLP to map each $x_i$ to a feature $h(x_i)$, applies a channel-wise global max pooling $g$, and a final MLP $\gamma$ for classification. To handle rigid transformations, PointNet includes a \textbf{T-Net} (Transformation Network) that learns a $3 \times 3$ alignment matrix for input coordinates and a $64 \times 64$ alignment matrix for feature spaces, enforced by a regularization term $\mathcal{R}(A) = \|I - A^T A\|_F$. For segmentation, the global feature vector is concatenated back to the individual point features before dense prediction.

PointNet's primary limitation is its failure to capture local geometric structures. \textbf{PointNet++} introduces a hierarchical stack of \textit{Set Abstraction Levels}:
1. \textbf{Sampling Layer}: Uses Farthest Point Sampling (FPS) to select well-distributed anchor points.
2. \textbf{Grouping Layer}: Defines local neighborhoods (e.g., metric balls) around anchors.
3. \textbf{PointNet Layer}: Extracts a local feature vector for each neighborhood.
To handle varying densities inherent in LiDAR, PointNet++ uses Multi-Scale Grouping (MSG) and Multi-Resolution Grouping (MRG) with random input dropout during training.

\subsection{Point Convolutional Operators (KPConv)}
KPConv bridges the gap between discrete convolutions and continuous point clouds. The convolution over a scattered neighborhood $\mathcal{N}_x$ is defined as:
\begin{equation}
(\mathcal{F} * g)(x) = \sum_{x_i \in \mathcal{N}_x} g(x_i - x) f_i
\end{equation}
The continuous kernel $g$ is parameterized by a set of $K$ learnable weight matrices $W_k$ associated with anchor points $\tilde{x}_k$:
\begin{equation}
g(y_i) = \sum_{k \le K} h(y_i, \tilde{x}_k) W_k
\end{equation}
where $h(y_i, \tilde{x}_k) = \max \left(0, 1 - \frac{\|y_i - \tilde{x}_k\|}{\sigma} \right)$ calculates the interpolation weights based on Euclidean distance.

\begin{insightbox}[Composite Convolutions]
Recent architectures propose separating spatial and feature processing. The Spatial Function extracts geometric structure, which is then multiplied with feature vectors. This factorization acts as a strong regularizer, allowing lighter network architectures to match KPConv's performance.
\end{insightbox}

\section{Autonomous Driving Applications}

Autonomous driving perception heavily relies on \textbf{3D Object Detection}. A 3D bounding box is parameterized by 7 values: $[x, y, z, h, w, l, yaw, class]$. Pitch and roll are usually neglected as vehicles adhere to the ground. LiDAR point clouds feature up to 150k points, 360-degree coverage, occlusions, and anisotropic densities. Real-time operation (>10 fps) is mandatory.

\subsection{Grid-Based Detectors: Voxels and Pillars}
\textbf{Bird Eye View (BEV)} is a highly effective 2D representation that preserves horizontal ranges while projecting the z-axis into channels, drastically reducing computational overhead.

\textbf{Voxel-based models} (e.g., \textbf{SECOND}) discretize the cloud into 3D voxels. To counteract extreme sparsity, \textbf{Submanifold Sparse Convolutions} are utilized: convolutions are only evaluated at active (non-empty) sites, preventing the spatial dilation of non-zero entries and eliminating computational overhead in empty space.

\textbf{PointPillars} takes this further by discretizing space only in the $x,y$ plane, creating vertical "pillars" of unlimited z-axis dimension. Each point is augmented to a 9-dimensional vector (coordinates, reflectance, offsets to cluster mean, offsets to pillar center). PointNet processes each pillar to yield a fixed-size representation. The scattered pillars are then mapped to a dense 2D BEV pseudo-image, seamlessly enabling rapid 2D CNN detection heads.

\subsection{Multimodal Learning: Fusion Strategies}
Autonomous vehicles simultaneously leverage LiDAR (precise 3D geometry but sparse) and RGB Cameras (dense semantics but lacking absolute depth). Calibrated extrinsic parameters permit spatial alignment.
\begin{itemize}
    \item \textbf{Early Fusion}: Intermediate features from both modalities are extracted, aligned, concatenated, and fused in a unified backbone. \textbf{BEVFusion} projects camera features into the BEV space using a View Transformer, concatenates them with LiDAR BEV features, and performs joint decoding.
    \item \textbf{Late Fusion}: Independent 3D and 2D detectors yield separate predictions, which are merged post-hoc via heuristic NMS or learned ensembles.
    \item \textbf{Cascade Fusion}: 2D camera detections generate 3D frustum proposals. LiDAR points inside these frustums are isolated to run targeted 3D localization. This is computationally demanding as it requires running the 3D network sequentially per proposal.
    \item \textbf{LCF3D (Late-Cascade Fusion)}: A state-of-the-art framework that executes late fusion first. The expensive cascade refinement is triggered \textit{only} for unmatched camera detections (candidate false negatives), dramatically improving recall for distant and small objects while maintaining real-time latency.
\end{itemize}
Stereo RGB-only or Monocular 3D detection architectures currently rank significantly lower on benchmarks (e.g., nuScenes, KITTI) than LiDAR or LiDAR+Camera systems, confirming that precise geometric measurements remain indispensable for reliable autonomous navigation.
